\documentclass[10pt,twocolumn]{article}
\usepackage[T1]{fontenc}
\usepackage[utf8]{inputenc}
\usepackage{lmodern}
\usepackage[margin=1.8cm,top=2.4cm,bottom=2cm,columnsep=0.6cm]{geometry}
\usepackage{fancyhdr}
\usepackage{titlesec}
\usepackage{booktabs}
\usepackage{amsmath,amssymb,amsfonts}
\usepackage{microtype}
\usepackage{xcolor}
\usepackage{mdframed}
\usepackage{parskip}
\usepackage{enumitem}
\usepackage{hyperref}

\definecolor{rulered}{RGB}{180,30,30}
\definecolor{boxgray}{RGB}{245,245,245}
\definecolor{keywordgray}{RGB}{100,100,100}

\pagestyle{fancy}
\fancyhf{}
\fancyhead[L]{\textsc{\small Claude Mag}}
\fancyhead[C]{\small\textit{Machine Learning Research Digest}}
\fancyhead[R]{\small 2026-08-17}
\fancyfoot[C]{\small\thepage}
\renewcommand{\headrulewidth}{0.8pt}

\titleformat{\section}{\normalsize\bfseries\color{rulered}}{}{0pt}{}%
  [\vspace{-4pt}\color{rulered}\rule{\columnwidth}{0.4pt}]
\titlespacing{\section}{0pt}{8pt}{4pt}

\hypersetup{colorlinks=true,urlcolor=rulered,linkcolor=black}

\begin{document}
\setlength{\parindent}{0pt}
\setlength{\parskip}{4pt}

{\small\color{keywordgray}\textbf{Keywords:}
  test-time scaling \textbullet{}
  retrieval-augmented reasoning \textbullet{}
  large reasoning models \textbullet{}
  chain-of-thought}\par\vspace{4pt}

{\LARGE\bfseries\raggedright
  ThinkRetrieve: Retrieval-Augmented Reasoning
  Traces for Test-Time Scaling}\par\vspace{4pt}

{\large\itshape\raggedright
  Dynamically Injecting Solved Exemplars Mid-Reasoning Counters Error Compounding
  and Improves Accuracy Across Five Models and Four Benchmarks}\par\vspace{6pt}

{\small\textbf{By} Vaibhav Singh, Soumya Suvra Ghosal, Sarvesh Gharat,
  Soumyabrata Pal, Ramasuri Narayanam, Dinesh Manocha}\par\vspace{2pt}

{\small\color{keywordgray}arXiv:2608.10928 \textbullet{} Preprint, August 2026 \textbullet{}
  IIT Bombay; Univ.\ of Maryland; Adobe Research}
\par\vspace{2pt}\noindent{\color{rulered}\rule{\columnwidth}{0.6pt}}\vspace{6pt}

Sequential test-time scaling of Large Reasoning Models (LRMs) often yields
diminishing returns: longer reasoning traces compound errors and drift from
the original problem. ThinkRetrieve fixes this by dynamically retrieving
solved step-by-step examples from an external corpus at each intermediate reasoning
step, injecting guidance on \emph{how} to reason rather than \emph{what facts}
are relevant. Experiments across five models (1.5B--8B parameters) on GSM-8K,
MATH-500, AIME 2025, and SciQ show consistent accuracy gains over standard
test-time scaling.

\begin{mdframed}[backgroundcolor=boxgray,linecolor=rulered,linewidth=1pt,
    innerleftmargin=6pt,innerrightmargin=6pt,
    innertopmargin=6pt,innerbottommargin=6pt]
\textbf{\small AT A GLANCE}\par\vspace{4pt}
\begin{description}[leftmargin=0pt,labelsep=4pt,itemsep=2pt,parsep=0pt,
    font=\small\bfseries]
  \item[Problem:] Allocating more inference compute (longer CoT traces) often
    backfires due to error compounding and problem drift in Large Reasoning Models.
  \item[Why it matters:] Test-time scaling is the dominant paradigm for improving LRM
    accuracy; fixing its failure mode is high-leverage and training-free.
  \item[Method:] At each reasoning step, retrieve the top-$m$ solved exemplars from an
    external corpus (conditioned on the partial trace) and inject them into the thinking trace.
  \item[Result:] Consistent accuracy gains on GSM-8K, MATH-500, AIME 2025, and SciQ
    across five models; largest gains on hard benchmarks and small models.
  \item[Evidence:] All five models (1.5B--8B) improve; gains replicated across four
    benchmarks spanning math and science reasoning.
\end{description}
\end{mdframed}

\section{The Big Picture}

Large Reasoning Models (LRMs) extend the chain-of-thought paradigm to its logical
limit: instead of a few reasoning steps, they generate extended thinking traces of
hundreds or thousands of tokens before producing a final answer. The more inference
compute allocated, the better the performance — in principle. In practice, sequential
test-time scaling saturates and then degrades: beyond some trace length, accuracy
falls because the model loses track of the problem, makes errors it cannot catch, and
compounds small mistakes into large final errors.

ThinkRetrieve reframes test-time scaling as a \emph{guided} process. Rather than
letting the model reason in isolation, it provides the model with solved analogues
— problems whose solutions are structurally similar to the current reasoning step —
at each point where the trace could diverge. The key insight is that retrieval is
most useful when it is \emph{dynamic}: the right exemplar at step 3 of a reasoning
trace is different from the right exemplar at step 10.

\section{Why Existing Approaches Fall Short}

Standard RAG retrieves documents before generation and injects them once into the
context. For reasoning tasks, this static approach misses the key failure mode:
errors do not arise from factual ignorance but from \emph{procedural} mistakes —
applying the wrong operation, misidentifying the relevant sub-problem, or making
an algebraic slip mid-trace. A retrieved document about a related topic does not
provide procedural guidance.

Best-of-$N$ sampling addresses error compounding by generating $N$ independent traces
and selecting the best one, but is expensive (linear in $N$) and does not actually
fix any individual trace — it relies on at least one trace avoiding the failure mode.
ThinkRetrieve fixes the trace \emph{at the point of failure}.

\section{Core Mechanism}

ThinkRetrieve wraps the LRM's generation loop with a retrieval step at each
reasoning checkpoint:

\begin{enumerate}[itemsep=2pt,parsep=0pt]
  \item At reasoning step $k$, form a query $q_k = [\text{problem}; \text{partial trace}_k]$.
  \item Retrieve top-$m$ solved exemplars $(p_j, \mathrm{sol}_j)$ from corpus $\mathcal{C}$
    based on embedding similarity to $q_k$.
  \item Inject the exemplars as additional in-context examples preceding step $k$.
  \item Continue the LRM's generation with the augmented context.
\end{enumerate}

The corpus $\mathcal{C}$ is a collection of (problem, step-by-step solution) pairs;
no fine-tuning of the LRM is performed. The retrieval encoder is a standard
dense retriever (e.g., a contrastively trained bi-encoder).

\section{Technical Formulation}

Let $\pi_\theta$ be the LRM policy and $x$ the input problem. Standard decoding
generates the trace $t = (t_1, \ldots, t_L)$ autoregressively:
\[
p(t \mid x) = \prod_{k=1}^L p_\theta(t_k \mid x, t_{<k})
\]
ThinkRetrieve modifies the conditioning at each step by inserting retrieved
exemplars $\mathcal{E}_k = \mathrm{Retrieve}(\mathcal{C},\, [x; t_{<k}])$:
\[
p^{\mathrm{TR}}(t_k \mid x, t_{<k}) = p_\theta(t_k \mid \mathcal{E}_k, x, t_{<k})
\]
where $\mathcal{E}_k$ is updated at each reasoning checkpoint (every $\delta$ tokens).
The final answer is extracted from $t_L$ as in standard decoding.

Retrieval similarity is cosine distance in the embedding space:
\[
\mathrm{sim}(q_k, p_j) = \frac{\mathrm{enc}(q_k)^\top\, \mathrm{enc}(p_j)}{\|\mathrm{enc}(q_k)\| \cdot \|\mathrm{enc}(p_j)\|}
\]
Top-$m$ exemplars by similarity are selected and prepended to the context.

\section{Learning or Inference Procedure}

No training is involved. The LRM ($\pi_\theta$) and dense retriever are both
frozen. The only hyperparameters are:
\begin{itemize}[itemsep=2pt,parsep=0pt]
  \item $m$: number of retrieved exemplars per step (ablated: 1, 3, 5)
  \item $\delta$: retrieval frequency (every $\delta$ generated tokens)
  \item Retrieval encoder: the paper uses a standard off-the-shelf bi-encoder
\end{itemize}

\section{What the Guarantee Says}

No formal guarantee is provided. The empirical claim is that dynamic step-conditioned
retrieval consistently improves accuracy over static retrieval and over no retrieval,
across model scales from 1.5B to 8B parameters and across four diverse benchmarks.
The consistency across scales and tasks provides strong evidence that the mechanism
is general rather than benchmark-specific.

\section{Experimental Findings}

\begin{itemize}[itemsep=2pt,parsep=0pt]
  \item \textbf{GSM-8K and MATH-500:} Consistent accuracy gains across all five models;
    gains increase with trace length (longer traces benefit most).
  \item \textbf{AIME 2025:} Largest absolute gains, where baseline traces are most
    error-prone and problem drift is most severe.
  \item \textbf{SciQ:} Gains confirmed on a non-math reasoning benchmark, establishing
    cross-domain generality.
  \item \textbf{Small models (1.5B--3B):} Largest relative gains; ThinkRetrieve provides
    the most leverage where baseline reasoning is weakest.
  \item \textbf{Ablation on $m$:} $m=3$ is near-optimal; $m=1$ is 60--80\% of the gain,
    $m>5$ provides marginal additional benefit at higher context cost.
\end{itemize}

\section{Ablations and Interpretation}

\textbf{Dynamic vs.\ static retrieval:} Retrieving once at the start (standard RAG)
provides $\sim$30\% of the accuracy gain of dynamic per-step retrieval — confirming
that the value is in the conditioning on the \emph{partial trace}, not just the problem.

\textbf{Retrieval timing:} Injecting exemplars early in the trace (step 1--2) is most
helpful for math benchmarks, where problem setup is the hardest phase; mid-trace
retrieval dominates for multi-step science problems.

\textbf{Corpus quality:} Using a corpus of correct solutions consistently outperforms
a corpus of incorrect-but-plausible solutions, confirming that the model is learning
procedural structure rather than just surface similarity.

\section*{Reference}

Vaibhav Singh, Soumya Suvra Ghosal, Sarvesh Gharat, Soumyabrata Pal,
Ramasuri Narayanam, Dinesh Manocha.
\textbf{ThinkRetrieve: Retrieval-Augmented Reasoning Traces for Test-Time Scaling}.
arXiv:2608.10928, August 2026.
\url{https://arxiv.org/abs/2608.10928}

\end{document}
